\documentclass[11pt]{article}
\usepackage[margin=1in]{geometry}
\usepackage{amsmath,amssymb}
\usepackage{hyperref}
\usepackage{booktabs}
\usepackage{graphicx}
\usepackage{svg}
\usepackage[numbers]{natbib}

\title{CoxMock Validation of the kNN Landy--Szalay Estimator}
\author{T.\ Abel \textit{et al.}\\ KIPAC / Stanford University}
\date{Draft --- March 2026}

\begin{document}
\maketitle

\begin{abstract}
We validate the $k$-Nearest Neighbor (kNN) pair-count-density estimator for
the two-point correlation function $\xi(r)$ using the CoxMock suite: an
isotropic line-point Cox process with analytically known $\xi(r)$.  Over
$N_{\rm mock} = 20$ realizations of the ``validation'' preset ($N_p = 10^4$,
$N_L = 10^3$, $\ell = 200\,h^{-1}\mathrm{Mpc}$, $L = 500\,h^{-1}\mathrm{Mpc}$),
the Davis--Peebles estimator recovers the analytic $\xi(r)$ with
$\chi^2/\mathrm{dof} = 1.22$ (30~bins, $r \in [3, 48]\,h^{-1}\mathrm{Mpc}$)
and maximum single-bin bias $\lvert\Delta\xi\rvert/\sigma_{\bar\xi} = 3.4$.
The kNN-CDF of the Poisson random catalog matches the Erlang prediction to
sub-percent level, confirming correct distance handling and tree queries.
\end{abstract}

% ------------------------------------------------------------------
\section{The CoxMock Generator}

We use an isotropic line-point Cox process following the Euclid 2PCF-GC
validation suite~\citep{delaTorre2025}.  $N_L$ lines of length~$\ell$ are
placed at random positions with isotropic random orientations
($\cos\theta$ uniform on $[-1,1]$, $\phi$ uniform on $[0,2\pi)$) in a
periodic cube of side~$L$.  Each of $N_p$ points is assigned to a
uniformly random line and placed at a uniformly random position~$t \in
[0,\ell]$ along that line.  Points whose coordinates exceed the box
boundaries are wrapped periodically into $[0, L)$.

\subsection{Analytic $\xi(r)$}

Each ordered pair of points has probability $1/N_L$ of sharing the same
line.  For same-line pairs the separation~$d = |t_1 - t_2|$ has the
triangular PDF~$f(d) = (2/\ell)(1 - d/\ell)$ for $0 \le d \le \ell$.
The expected number of same-line ordered pairs is
$N_p(N_p - 1)/N_L$.  Dividing the excess in a spherical shell $4\pi
r^2\,\mathrm{d}r$ by the Poisson expectation $\bar{n}\,4\pi
r^2\,\mathrm{d}r$ yields:
\begin{equation}
  \xi(r)
    = \frac{N_p - 1}{N_L\,\bar{n}\,2\pi\,r^2\,\ell}
      \;\Bigl(1 - \frac{r}{\ell}\Bigr)\,,
  \qquad 0 < r < \ell\,,
  \label{eq:xi_cox}
\end{equation}
with $\xi(r) = 0$ for $r \ge \ell$.  Here $\bar{n} = N_p/V$,
$V = L^3$.  For $N_p \gg 1$ this simplifies to $\xi(r) \approx V/(N_L\,
2\pi\,r^2\,\ell)\,(1 - r/\ell)$.

\textbf{Note:} The prefactor uses $N_p - 1 \approx m\,N_L$ (not
$m(m{-}1)\,N_L$) because the generator assigns points to lines
\emph{randomly}, giving Multinomial line-occupation numbers
$m_i \sim \mathrm{Multinomial}(N_p; 1/N_L, \ldots, 1/N_L)$ with
$E[\sum m_i(m_i{-}1)] = N_p(N_p{-}1)/N_L$.  For fixed-$m$ assignment
the prefactor would be $(m{-}1)/(\bar{n}\,2\pi\,r^2\,\ell)$ instead.

The profile is $\xi \propto r^{-2}\,(1 - r/\ell)$, so the pair-count
contribution $r^2\,\xi(r) \propto (1 - r/\ell)$ is nearly flat at small
$r$ and falls linearly to zero at $r = \ell$.

\paragraph{Validation preset parameters.}
\begin{center}
\begin{tabular}{@{}lrl@{}}
\toprule
Parameter & Value & Description \\
\midrule
$L$           & $500\,h^{-1}\mathrm{Mpc}$ & Box side length \\
$N_L$         & 1\,000   & Number of lines \\
$\ell$        & $200\,h^{-1}\mathrm{Mpc}$ & Line length \\
$N_p$         & 10\,000  & Number of data points \\
$m = N_p/N_L$ & 10       & Mean points per line \\
$\bar{n}$     & $8 \times 10^{-5}\,h^3\mathrm{Mpc}^{-3}$ & Mean density \\
\bottomrule
\end{tabular}
\end{center}

% ------------------------------------------------------------------
\section{Estimator}

The pair-count density $\mathcal{D}(r)$ is obtained by histogramming
the $k_{\max}$ nearest-neighbor distances from a set of query points
into radial bins $[r_i, r_{i+1})$ and normalizing by the number of
queries and the shell volume $4\pi r_{\rm mid}^2\,\Delta r$:
\begin{equation}
  \mathcal{D}(r_i)
    = \frac{1}{N_{\rm query}\,4\pi\,r_i^2\,\Delta r}
      \sum_{q=1}^{N_{\rm query}} \#\{k : d_{q,k} \in [r_i, r_{i+1})\}\,.
\end{equation}

For the data--data term, each data point queries the data tree for
$k_{\max} + 1$ neighbors; the self-pair ($d \approx 0$) is filtered,
leaving $k_{\max}$ real neighbors.  For data--random (DR), $N_R =
10\,N_p$ uniform random points query the data tree for $k_{\max}$
neighbors each.

The Davis--Peebles estimator for a periodic box (where $\mathcal{D}^{\rm
RR} = \bar{n}$ is constant) reduces to:
\begin{equation}
  \hat\xi(r) = \frac{\mathcal{D}^{\rm DD}(r)}{\mathcal{D}^{\rm DR}(r)} - 1\,.
  \label{eq:dp_knn}
\end{equation}

All distance computations use periodic (minimum-image) boundary
conditions matching the CoxMock generator.

\paragraph{Run parameters.}
$k_{\max} = 128$ ($r_{k_{\max}} \approx 72.6\,h^{-1}\mathrm{Mpc}$,
well beyond $r_{\max} = 48$); $N_R/N_d = 10$; 30~linear bins in $r
\in [3, 48]\,h^{-1}\mathrm{Mpc}$; 20~independent mock realizations.

% ------------------------------------------------------------------
\section{Results}

\subsection{$\xi(r)$ Recovery}

Figure~\ref{fig:xi_comparison} shows the mean $\hat\xi(r)$ over 20
mocks (blue line with $\pm 1\sigma$ band) compared to the analytic
prediction (Eq.~\ref{eq:xi_cox}, black dashed).  The agreement is
excellent across the full range $r \in [3, 48]\,h^{-1}\mathrm{Mpc}$.
The signal spans nearly three orders of magnitude, from
$\xi(3.8) \approx 6.9$ down to $\xi(47) \approx 0.03$.

\begin{figure}[ht]
  \centering
  \includesvg[width=0.9\textwidth]{xi_vs_analytic}
  \caption{Mean $\hat\xi(r)$ from 20 CoxMock realizations (blue, with
  $\pm 1\sigma$ band) vs.\ the analytic
  $\xi(r) = (N_p{-}1)/(N_L\,\bar{n}\,2\pi r^2\ell)\,(1-r/\ell)$
  (black dashed).}
  \label{fig:xi_comparison}
\end{figure}

\subsection{$r^2\,\xi(r)$ Representation}

The $r^2\,\xi(r)$ representation (Fig.~\ref{fig:r2xi}) removes the
$1/r^2$ divergence, revealing the nearly flat plateau at small~$r$ and
the linear drop to zero at $r = \ell$.  The measured curve closely
follows the analytic $(1 - r/\ell)$ profile.

\begin{figure}[ht]
  \centering
  \includesvg[width=0.9\textwidth]{r2xi_comparison}
  \caption{$r^2\,\hat\xi(r)$ (blue, with $\pm 1\sigma$ band) vs.\ the
  analytic $r^2\,\xi(r) \propto (1 - r/\ell)$ (black dashed).}
  \label{fig:r2xi}
\end{figure}

\subsection{Bias and $\chi^2$}

Figure~\ref{fig:residuals} shows the residuals
$(\bar\xi - \xi_{\rm true})/\sigma_{\bar\xi}$ in each bin.
All residuals lie within $\pm 3.4\sigma$, consistent with the expected
scatter for 30 bins.  The maximum single-bin deviation is $3.38\sigma$
at $r = 38.2\,h^{-1}\mathrm{Mpc}$.

\begin{figure}[ht]
  \centering
  \includesvg[width=0.9\textwidth]{xi_residuals}
  \caption{Residuals $(\bar\xi - \xi_{\rm true})/\sigma_{\bar\xi}$ in
  each radial bin.  The $\pm 2\sigma$ bands (crimson dashed) are shown
  for reference.}
  \label{fig:residuals}
\end{figure}

\begin{table}[ht]
  \centering
  \caption{Summary statistics for the CoxMock validation.}
  \label{tab:summary}
  \begin{tabular}{@{}lc@{}}
    \toprule
    Quantity & Value \\
    \midrule
    $N_{\rm mock}$ & 20 \\
    Radial bins & 30 linear in $[3, 48]\,h^{-1}\mathrm{Mpc}$ \\
    $k_{\max}$ & 128 \\
    $N_R / N_d$ & 10 \\
    $\chi^2/\mathrm{dof}$ & $36.5/30 = 1.22$ \\
    Max $|\mathrm{bias}|/\sigma_{\bar\xi}$ & 3.38 \\
    $p(\chi^2 \ge 36.5 \mid 30\,\mathrm{dof})$ & $\approx 0.19$ \\
    \bottomrule
  \end{tabular}
\end{table}

\subsection{kNN-CDF vs.\ Erlang}

For the uniform random catalog, the kNN cumulative distribution
functions must match the Erlang distribution:
\begin{equation}
  \mathrm{CDF}_k(r)
    = 1 - e^{-\lambda}\sum_{j=0}^{k-1}\frac{\lambda^j}{j!}\,,
  \qquad \lambda = \bar{n}_R\,\tfrac{4}{3}\pi r^3\,.
\end{equation}
Figure~\ref{fig:cdf_erlang} confirms sub-percent agreement for
$k = 1, 2, 4, 8$, validating the tree queries and periodic distance
computation.

\begin{figure}[ht]
  \centering
  \includesvg[width=0.9\textwidth]{cdf_comparison}
  \caption{Measured kNN-CDFs for the random catalog (colored lines) vs.\
  Erlang predictions (black dashed) for $k = 1, 2, 4, 8$.}
  \label{fig:cdf_erlang}
\end{figure}

\subsection{Mock-to-Mock Scatter}

Figure~\ref{fig:individual} overlays all 20 individual mock
$\hat\xi(r)$ curves (light blue) along with the mean (dark blue) and
the analytic prediction (black dashed).  The mock-to-mock scatter is
visually consistent with Gaussian fluctuations around the mean, with no
outlier realizations.

\begin{figure}[ht]
  \centering
  \includesvg[width=0.9\textwidth]{individual_mocks}
  \caption{All 20 individual $\hat\xi(r)$ curves (light blue), their
  mean (bold blue), and the analytic $\xi(r)$ (black dashed).}
  \label{fig:individual}
\end{figure}

\subsection{Combined Summary}

Figure~\ref{fig:summary} presents a combined 2$\times$2 panel showing
(a)~$\xi(r)$, (b)~$r^2\xi(r)$, (c)~residuals, and (d)~kNN-CDF
comparison.

\begin{figure}[ht]
  \centering
  \includesvg[width=\textwidth]{validation_summary}
  \caption{Combined validation summary: (a) $\xi(r)$ recovery,
  (b) $r^2\xi(r)$, (c) bias residuals, (d) kNN-CDF vs.\ Erlang.}
  \label{fig:summary}
\end{figure}

% ------------------------------------------------------------------
\section{Discussion}

The kNN pair-count-density Davis--Peebles estimator recovers the
analytic $\xi(r)$ of the isotropic line-point Cox process to within the
statistical precision of $N_{\rm mock} = 20$ realizations.  Key
findings:

\begin{enumerate}
\item \textbf{Unbiased recovery.}  The $\chi^2/\mathrm{dof} = 1.22$ is
  consistent with statistical expectations ($p \approx 0.19$).  All bins
  have $|\mathrm{bias}|/\sigma_{\bar\xi} < 3.4$.

\item \textbf{Correct analytic formula.}  For random line assignment
  (Multinomial occupation numbers), the prefactor involves
  $E[\sum m_i(m_i{-}1)] = N_p(N_p{-}1)/N_L \approx m\,N_p$, not
  $m(m{-}1)\,N_L$.  This differs from the deterministic-assignment
  formula by a factor $m/(m{-}1)$, which is 11\% for $m = 10$.

\item \textbf{Periodic boundary conditions.}  Using minimum-image
  periodic distances is essential; without them, boundary effects
  introduce subtle biases at all scales.

\item \textbf{$k_{\max}$ choice.}  With $k_{\max} = 128$, the
  characteristic $k$NN distance is $r_{k_{\max}} \approx 72.6\,
  h^{-1}\mathrm{Mpc}$, safely beyond $r_{\max} = 48$.  Smaller values
  ($k_{\max} = 32$ or 64) produce truncation artifacts: clustered data
  points use many kNN slots for nearby same-line neighbors, depleting
  counts at large~$r$ and causing a systematic negative bias in the
  outer bins.

\item \textbf{Self-pair exclusion.}  DD queries use $k_{\max} + 1$
  neighbors; the self-match at $d \approx 0$ is filtered out, ensuring
  exactly $k_{\max}$ real neighbors per query.

\item \textbf{CDF sanity check.}  The kNN-CDF of the Poisson random
  catalog matches the Erlang distribution, confirming that the brute-force
  $k$NN implementation and periodic distance calculation are correct.
\end{enumerate}

This validation establishes the correctness of the kNN pair-count-density
estimator for a point process with $\xi$ spanning $\sim 3$ orders of
magnitude ($\xi \sim 0.03$--$7$).  The next steps are: (i)~extension to
non-periodic geometries with survey masks, (ii)~implementation of
the dilution-ladder for probing multiple density regimes, and
(iii)~comparison with the standard Landy--Szalay pair-count estimator on
the same mocks.

\end{document}
